% !TeX root = ./main.tex

% --------------------------------------------------
% 資訊設定（Information Configs）
% --------------------------------------------------

\ntusetup{
  university*   = {National Taiwan University},
  university    = {國立臺灣大學},
  college       = {電機資訊學院},
  college*      = {College of Electronical Engineering and Computer Science},
  institute     = {資訊工程研究所},
  institute*    = {Department of Computer Science an Information Engineering}, % 請確認您的科系為Department of xxx或Institute of xxx
  title         = {應用於邊緣運算環境之低延遲 Wasm 執行時期分層編譯架構},
  title*        = {A Tiered Compilation Architecture for Low-Latency Wasm Runtimes in Edge
Computing Environments},
  author        = {李祥宇},
  author*       = {Hsiang-Yu Lee},
  ID            = {R12922137},
  advisor       = {廖世偉},
  advisor*      = {Shih-Wei Liao},
  %date          = {2020-05-01},         % 若註解掉，則預設為當天
  %oral-date     = {2020-05-01},         % 若註解掉，則預設為當天
  DOI           = {10.5566/NTU2018XXXXX},
  keywords      = {WebAssembly, 即時編譯, 分層編譯, 在堆疊替換, 邊緣運算, WasmEdge},
  keywords*     = {WebAssembly, just-in-time compilation, tiered compilation, on-stack replacement, edge computing, WasmEdge},
}

% --------------------------------------------------
% 加載套件（Include Packages）
% --------------------------------------------------

\usepackage[sort&compress]{natbib}      % 參考文獻
\usepackage{amsmath, amsthm, amssymb}   % 數學環境
\usepackage{ulem, CJKulem}              % 下劃線、雙下劃線與波浪紋效果
\usepackage{booktabs}                   % 改善表格設置
\usepackage{multirow}                   % 合併儲存格
\usepackage{diagbox}                    % 插入表格反斜線
\usepackage{array}                      % 調整表格高度
\usepackage{longtable}                  % 支援跨頁長表格
\usepackage{paralist}                   % 列表環境


\usepackage{lipsum}                     % 英文亂字
\usepackage{zhlipsum}                   % 中文亂字

\usetikzlibrary{positioning}            % tikz: relative-position syntax (below=of, etc.)

\usepackage{algorithm}                  % algorithm float environment
\usepackage{algpseudocode}              % algorithmic pseudocode (\State, \If, \While, ...)
\usepackage{listings}                   % verbatim code listings with captions
\lstset{
  basicstyle=\ttfamily\footnotesize,
  breaklines=true,
  captionpos=b,
  frame=single,
  showstringspaces=false,
}

% --------------------------------------------------
% TOC and section-numbering overrides
% --------------------------------------------------

% TOC includes chapters, sections, and subsections.
\setcounter{tocdepth}{2}
\setcounter{secnumdepth}{2}

% A package in the load chain (likely titletoc / tocloft) bypasses tocdepth
% for \paragraph and adds every \paragraph heading to the TOC regardless of
% the tocdepth setting. Redefine \paragraph to invoke its starred form
% unconditionally: same visual appearance (bold run-in head), no TOC entry.
\let\origparagraph\paragraph
\renewcommand{\paragraph}[1]{\origparagraph*{#1}}

% --------------------------------------------------
% Appendix command override
% --------------------------------------------------

% The class's \appendix renames \thechapter to the appendix letter but does
% not increment the underlying chapter counter, so each appendix's section
% anchors collide (every appendix's first section has anchor section.N.1
% for the same N). Override to refstep the chapter counter so each appendix
% has a unique anchor namespace, while still printing the appendix letter
% as the chapter label.
\makeatletter
\renewcommand{\appendix}[2]{
  \refstepcounter{chapter}
  \chapter*{\ntu@appendix~{#1}~---~#2}
  \phantomsection
  \addcontentsline{toc}{chapter}{\ntu@appendix~{#1}~---~#2}
  \setcounter{section}{0}
  \renewcommand{\thechapter}{#1}
}
\makeatother

% --------------------------------------------------
% 套件設定（Packages Settings）
% --------------------------------------------------
